\documentclass[11pt]{article}
\usepackage{amsmath,amssymb,amsthm,mathtools}
\usepackage[margin=1in]{geometry}
\usepackage[hidelinks]{hyperref}

\newtheorem{theorem}{Theorem}
\newtheorem{proposition}[theorem]{Proposition}
\newtheorem{corollary}[theorem]{Corollary}
\newtheorem{definition}[theorem]{Definition}
\newtheorem{remark}[theorem]{Remark}

\newcommand{\dd}{\mathrm{d}}
\DeclareMathOperator{\sgn}{sgn}

\title{Parity--fixed inverse--channel normal forms\\ and universal curvature blow--up}
\author{[authors]}
\date{Draft --- \today}

\begin{document}
\maketitle

\begin{abstract}
We isolate the local curvature calculus behind the complementarity law of a diagonal
inverse--channel thermodynamic metric. Near a boundary where one output coordinate collapses
against transverse channels, the scalar curvature has a universal Laurent normal form
determined entirely by the local germ and its parity. A \emph{generic} quadratic collapse with
finite transverse channels gives an order--$3$ pole whose coefficient is the transverse first
drift; a \emph{parity--fixed} collapse, where the transverse channels themselves diverge and the
germ is even, gives an order--$4$ pole with the dimension--universal coefficient
$-m(m+5)/B$, where $m$ is the transverse fibre dimension and $B$ the collapsing amplitude---the
transverse amplitudes are forgotten. For a three--dimensional state space ($m=2$) this is
$-14/B$. Kerr--Newman thermodynamic geometry is the $m=2$ application: the extremal edge is a
generic face (order $3$), the two charge--reflection faces are parity--fixed (order $4$,
$-14/B$), and their corner is a higher--codimension blow--up stratum. All statements are
checked by a reproducible symbolic certificate.
\end{abstract}

%======================================================================
\section{The local setting}

Let $g=\operatorname{diag}(g_x,g_{y_1},\dots,g_{y_m})$ be a diagonal metric on an
$(m{+}1)$--dimensional state space near a face $\{x=0\}$ across which the native coordinate $x$
collapses. Write $R$ for the scalar curvature. Two germs occur, distinguished by whether the
transverse channels stay finite or themselves diverge, and by parity in $x$.

\paragraph{Convention.} A germ is \emph{parity--fixed} if every component is even in $x$
(invariant under $x\mapsto-x$); then $R$ is even in $x$ and all odd--order poles vanish
identically.

%======================================================================
\section{The pure parity--fixed normal form}

\begin{theorem}[Pure parity--fixed normal form]
\label{thm:pure}
Let $m\ge1$, $B>0$, $A_\alpha>0$. For
\[
\dd s^2=Bx^2\,\dd x^2+x^{-2}\sum_{\alpha=1}^{m}A_\alpha\,\dd y_\alpha^2,
\]
the scalar curvature is exactly
\[
R=-\frac{m(m+5)}{B}\,x^{-4},
\]
independent of all transverse amplitudes $A_\alpha$. In particular $m=2$ gives $R=-14\,x^{-4}/B$.
\end{theorem}

\begin{proof}
Set $r=\tfrac12\sqrt B\,x^2$, so $\dd r=\sqrt B\,x\,\dd x$ and $x^{-2}=\sqrt B/(2r)$. The metric
is the warped product $\dd s^2=\dd r^2+w(r)^2 h_{\mathrm{flat}}$ with flat fibre
$h_{\mathrm{flat}}=\sum_\alpha A_\alpha\,\dd y_\alpha^2$ and $w^2=\sqrt B/(2r)$, i.e.\ $w\propto
r^{-1/2}$. For a warped product over a flat $m$--dimensional fibre,
$R=-2m\,w''/w-m(m-1)(w'/w)^2$; with $w'/w=-1/(2r)$ and $w''/w=3/(4r^2)$,
\[
R=-\frac{6m}{4r^2}-\frac{m(m-1)}{4r^2}=-\frac{m(m+5)}{4r^2},
\]
and $4r^2=Bx^4$. The two contributions are radial focusing ($6m$) and transverse fibre shear
($m(m-1)$); the amplitudes cancel because they only rescale a flat fibre, and warp curvature
sees only $w'/w,w''/w$.
\end{proof}

\begin{corollary}[Asymptotic parity--fixed face]
\label{cor:asymp}
If near $x=0$, with all components even in $x$,
\[
g_x=B(y)x^2+O(x^4),\qquad g_{y_\alpha}=A_\alpha(y)x^{-2}+O(1),\qquad B,A_\alpha>0,
\]
then $R=-m(m+5)\,x^{-4}/B(y)+O(x^{-2})$. The odd pole $x^{-3}$ is absent by parity, and the
leading coefficient is independent of the amplitudes $A_\alpha(y)$ and of all even subleading
data.
\end{corollary}

%======================================================================
\section{The generic collapse, and why order jumps by one}

\begin{proposition}[Generic quadratic collapse]
\label{prop:generic}
Let the transverse channels stay finite and the germ break parity at first order:
\[
g_x=Ax^2+O(x^3),\qquad g_{y_\alpha}=P_{\alpha0}+P_{\alpha1}x+O(x^2),\qquad A,P_{\alpha0}\neq0.
\]
Then
\[
R=\frac1A\sum_{\alpha=1}^{m}\frac{P_{\alpha1}}{P_{\alpha0}}\,x^{-3}+O(x^{-2})
=\frac1A\sum_{\alpha}\partial_x\log P_\alpha\Big|_{0}\,x^{-3}+O(x^{-2}),
\]
independent of $A_3$ and of the quadratic transverse data $P_{\alpha2}$.
\end{proposition}

The two propositions sit side by side and expose the mechanism. A generic face is order $3$
because the odd pole is produced by the parity--\emph{breaking} linear drift $P_{\alpha1}$ of
the transverse channels. A parity--fixed face has no such drift---the germ is even---so the
$x^{-3}$ term vanishes identically and the leading pole is the even $x^{-4}$ term, whose
coefficient is the dimension--universal $-m(m+5)/B$. The order is not an accident of a
particular fundamental relation; it is read off the parity of the germ.

%======================================================================
\section{Higher--codimension corners}
\label{sec:corner}

Where several parity variables vanish together the leading order is direction--dependent. On a
two--variable corner with $x=\epsilon$ and a second collapsing coordinate scaled as
$\epsilon^{a}$, the curvature order is piecewise linear in $a$: a Newton polygon whose vertices
are the balanced directions and whose edges interpolate to the single--face orders. The order
can \emph{drop} below either face---two order--$4$ faces can meet at a corner whose balanced
approach is only order $2$---and a projective front--face coefficient records the
path--dependence that neither single face sees. The general corner theorem (a valuation method
that reads the order off the germ without a global scalar expansion) is left to a companion
note; the Kerr--Newman Schwarzschild corner below is the worked instance.

%======================================================================
\section{Kerr--Newman as the $m=2$ application}

For the Kerr--Newman thermodynamic surface the inverse--channel DBP metric is diagonal in the
role coordinates $(S,J,Q)$, a three--dimensional state space, so $m=2$. The local types are:

\begin{center}
\begin{tabular}{llcl}
\hline
divisor & local type & order & leading coefficient\\
\hline
$T=0$ (extremal) & generic collapse (Prop.~\ref{prop:generic}) & $3$ &
  $A_2^{-1}\,\partial_S\log(G_JG_Q)|_{\mathrm{ext}}$\\
$\Omega=0$ (spin) & parity--fixed (Cor.~\ref{cor:asymp}) & $4$ & $-14/B_J$\\
$\Phi_e=0$ (charge) & parity--fixed (Cor.~\ref{cor:asymp}) & $4$ & $-14/B_Q$\\
$\Omega=\Phi_e=0$ (Schwarzschild) & double corner (\S\ref{sec:corner}) & wedge & mixed stratum\\
\hline
\end{tabular}
\end{center}

The extremal order--$3$ and its sign, the two reflection coefficients $-14/B$, and the
corner's order--$2$ balanced approach with its mixed excess are all established in the
Kerr--Newman study; here they are exhibited as instances of the local classification, not as
surface--specific miracles. The number $14$ is $6\cdot2+2\cdot1$: radial focusing plus
transverse fibre shear for a two--dimensional fibre.

%======================================================================
\section*{Certificate}

The pure normal form (Thm.~\ref{thm:pure}, first--principles Ricci for $m=1,\dots,4$ and the
symbolic--in--$m$ warped reduction), the asymptotic parity face (Cor.~\ref{cor:asymp}), and the
generic collapse (Prop.~\ref{prop:generic}) are machine--checked in
\texttt{verification/pfc\_test1\_local\_normal\_forms.py}, run twice with identical output. The
Kerr--Newman instances are in \texttt{lead7\_test6,7,8,9,10}.

\end{document}
